% ----------------------------------------------------------------------
%  3.4.2. Макроскопическое описание диффузии
% ----------------------------------------------------------------------
\subsection{Макроскопическое описание диффузии}\label{ssec:fick}

На макроуровне твёрдое тело рассматривают как сплошную среду, забывая
об атомах. Основа~— \term{уравнение непрерывности} (вещество не
возникает и не исчезает):
\begin{equation}\label{eq:continuity}
\frac{\partial C}{\partial t}+\operatorname{div}\vb J=0 .
\end{equation}

\paragraph{Законы Фика}
При отжиге неоднородного однофазного сплава возникает поток вещества,
выравнивающий концентрацию. \term{Первый закон Фика}:
\begin{keyeq}\label{eq:fick1}
J=-D\frac{\partial C}{\partial x}\qquad\text{или}\qquad \vb J=-D\,\nabla C ,
\end{keyeq}
где $D$~— \term{коэффициент диффузии}. Его размерность~— м$^2$/с, но в
справочниках чаще встречаются см$^2$/с ($1~\text{см}^2/\text{с}=10^{-4}~\text{м}^2/\text{с}$).
Подставив \eqref{eq:fick1} в \eqref{eq:continuity}, получим
\term{второй закон Фика}~— как меняется концентрация во времени:
\begin{keyeq}\label{eq:fick2}
\frac{\partial C}{\partial t}=\frac{\partial}{\partial x}\left(D\frac{\partial C}{\partial x}\right)
\quad\xrightarrow{\ D=\text{const}\ }\quad
\frac{\partial C}{\partial t}=D\frac{\partial^2C}{\partial x^2}.
\end{keyeq}
В трёхмерном случае при постоянном $D$: $\partial C/\partial t=D\nabla^2C$.
Фик получил эти уравнения по аналогии с теплопроводностью, и с тех пор
они не изменились.

\begin{table}[ht]
\centering
\caption{Порядки коэффициентов диффузии}
\label{tab:D-orders}
\small
\begin{tabular}{@{}lc@{}}
\toprule
Случай & $D$, м$^2$/с \\
\midrule
атомы в жидком металле & $\sim10^{-9}$ \\
углерод в $\gamma$-Fe (внедрение), 1000\,°C & $\sim10^{-11}$ \\
самодиффузия металлов у точки плавления & $10^{-12}$--$10^{-10}$ \\
самодиффузия меди при комнатной температуре & $\sim10^{-40}$ \\
\bottomrule
\end{tabular}
\end{table}

\begin{glavnoe}
Из второго закона Фика следует простая и очень полезная оценка:
за время $t$ диффузия распространяется на расстояние порядка
\begin{equation}\label{eq:diff-length}
L\sim\sqrt{Dt}.
\end{equation}
Расстояние растёт как \emph{корень} из времени: чтобы продвинуться вдвое
дальше, нужно в четыре раза больше времени. Например, при
$D=10^{-12}$~м$^2$/с за час ($3600$~с) получаем
$L\sim\sqrt{3{,}6\cdot10^{-9}}\approx60$~мкм.
\end{glavnoe}

\begin{masshtab}[диффузионный путь за 1 час]{-4.3/-3.3/{самодиффузия у $T_\text{пл}$}, -3.85/-3.55/{C в $\gamma$-Fe, 1000\,°C}, -2.85/-2.6/{жидкий металл}}
Подставим в $L\sim\sqrt{Dt}$ значения из табл.~\ref{tab:D-orders} и $t=1$~ч.
У точки плавления самодиффузия проходит 60--600~мкм~— больше размера
зерна, поэтому за часы отжига состав выравнивается во всём зерне.
Углерод в аустените при 1000\,°C уходит на $\sim0{,}2$~мм (на этом основана
цементация), атомы в жидкости~— на $\sim2$~мм. А самодиффузия меди при
комнатной температуре ($D\sim10^{-40}$~м$^2$/с) даёт $L\sim10^{-18}$~м~—
в миллион раз меньше размера атома: за час не происходит ни одного
перескока, и на линейке этот путь просто не помещается.
\end{masshtab}

\paragraph{Стационарные решения}
В стационарном состоянии $\partial C/\partial t=0$, и нужно решить
$\nabla^2C=0$. В одномерных случаях:
\[
\frac{\d^2C}{\d x^2}=0\ \ \text{(плоскость)},\qquad
\frac{\d^2C}{\d r^2}+\frac1r\frac{\d C}{\d r}=0\ \ \text{(цилиндр)},\qquad
\frac{\d^2C}{\d r^2}+\frac2r\frac{\d C}{\d r}=0\ \ \text{(сфера)}.
\]
Решения: $C$ линейна по $x$, по $\ln r$ и по $1/r$ соответственно.
Например, поток через пластину толщиной $h$ при концентрациях $C_1$ и
$C_2$ на её сторонах постоянен: $J=D(C_1-C_2)/h$. Так рассчитывают
проникновение водорода сквозь стенку сосуда или оболочку твэла.

\paragraph{Тонкий слой: гауссово решение}
Нанесём на торец образца тонкий слой диффузанта (его количество на
единицу площади~— $\alpha$) и отожжём в течение времени $t$. Если $D$
постоянен, решение второго закона Фика~— гауссиана:
\begin{keyeq}\label{eq:gauss}
C(x,t)=\frac{\alpha}{\sqrt{\pi Dt}}\exp\!\left(-\frac{x^2}{4Dt}\right),
\qquad \int_0^\infty C\,\d x=\alpha ,
\end{keyeq}
где $x$~— расстояние от торца\ispr{В оригинале (с.~430) множитель
равен $\alpha/(2\sqrt{\pi Dt})$, а условие сохранения записано как
$\int_{-\infty}^{\infty}C\,\d x=\alpha$. Это решение для тонкого слоя
\emph{внутри} бесконечного образца (между двумя половинами), когда
вещество уходит в обе стороны. Для слоя на торце полубесконечного
образца, как описано в тексте, вещество уходит только в одну сторону, и
множитель вдвое больше.}. Со временем профиль расплывается: высота
падает как $1/\sqrt t$, а ширина растёт как $\sqrt{Dt}$
(рис.~\ref{fig:gauss}).

\begin{figure}[ht]
\centering
\begin{tikzpicture}
\begin{axis}[width=0.62\textwidth, height=5cm, xmin=0, xmax=4, ymin=0, ymax=1.2,
  xlabel={$x$}, ylabel={$C$}, xtick=\empty, ytick=\empty, axis lines=left]
  \addplot[fmBlue, thick, domain=0:4, samples=100] {1/sqrt(pi*0.25)*exp(-x^2/(4*0.25))*0.6};
  \addplot[fmGreen!70!black, thick, domain=0:4, samples=100] {1/sqrt(pi*1)*exp(-x^2/4)*0.6};
  \addplot[fmRed, thick, domain=0:4, samples=100] {1/sqrt(pi*4)*exp(-x^2/16)*0.6};
  \node[font=\small, fmBlue] at (axis cs:0.75,0.62) {$t_1$};
  \node[font=\small, fmGreen!60!black] at (axis cs:1.6,0.3) {$4t_1$};
  \node[font=\small, fmRed] at (axis cs:3.0,0.18) {$16t_1$};
\end{axis}
\end{tikzpicture}
\caption{Расплывание тонкого слоя: при увеличении времени в 4 раза
ширина профиля удваивается, а высота уменьшается вдвое (площадь под
кривыми одинакова)}
\label{fig:gauss}
\end{figure}

\paragraph{Как измеряют коэффициент диффузии}
Из \eqref{eq:gauss} $\ln C=\text{const}-x^2/4Dt$: в координатах
$\ln C$~— $x^2$ профиль~— прямая с наклоном $-1/(4Dt)$. Значит, нужно
после отжига снимать тонкие слои параллельно торцу, измерять в каждом
концентрацию и по наклону прямой найти $D$. Почти всегда диффузант
берут в виде \emph{радиоактивного изотопа}: это резко повышает
чувствительность и, главное, позволяет измерять \emph{самодиффузию}~—
движение атомов металла в собственной решётке. Коэффициенты, измеренные
методом меченых атомов, обозначают $D^*$.

\begin{podrobnee}[метод Грузина и пределы измерений]
П.\,Л.~Грузин предложил измерять не активность снятого слоя, а
интегральную активность оставшейся части образца. Излучение глубже
лежащих слоёв ослабляется с коэффициентом поглощения $\mu$:
\[
I(x_n)=\int_{x_n}^\infty i(x)\,e^{-\mu(x-x_n)}\d x ,
\]
где $i(x)\propto C(x,t)$. Дифференцируя, получаем
$i(x_n)=\mu I_n-\d I_n/\d x_n\propto\exp(-x_n^2/4D^*t)$, и $D^*$
находят тем же способом.

Верхний предел измерения ($10^{-8}$--$10^{-6}$~см$^2$/с) соответствует
самодиффузии у точки плавления. Нижний предел определяется тем, что
концентрация в снятых слоях должна измениться хотя бы на порядок:
$x^2/4D^*t\approx2{,}3$, т.\,е. $x\approx3\sqrt{D^*t}$. При толщине слоя
около 0,1~мкм\isprmark{} ($x^2=10^{-10}$~см$^2$) и времени отжига $10^4$~ч
\[
D^*\approx\frac{x^2}{9t}=\frac{10^{-10}}{9\cdot10^4\cdot3600}\approx3\cdot10^{-19}~\text{см}^2/\text{с}.
\]
Для слоёв толщиной 1~мкм предел был бы около
$3\cdot10^{-17}$~см$^2$/с.
\end{podrobnee}
\isprtext{В оригинале (с.~432) говорится о слоях толщиной «более
микрона», но в расчёт подставлено $x^2=10^{-10}$~см$^2$, т.\,е.
$x=0{,}1$~мкм, и множитель 9 потерян; получено $\approx10^{-18}$~см$^2$/с.}

\paragraph{Контакт двух полубесконечных тел}
Пусть исходно $C=0$ при $x<0$ и $C=C_0$ при $x>0$ (рис.~\ref{fig:erf},\,\textit{а}).
Разобьём правую половину на тонкие слои толщиной $\Delta\alpha$; каждый
расплывается в гауссиану, а истинное распределение~— их сумма
(рис.~\ref{fig:erf},\,\textit{б}). В пределе сумма становится интегралом,
и решение выражается через \term{функцию ошибок}
$\operatorname{erf}z=\frac2{\sqrt\pi}\int_0^ze^{-\eta^2}\d\eta$:
\begin{keyeq}\label{eq:erf}
C(x,t)=\frac{C_0}2\left[1+\operatorname{erf}\frac{x}{2\sqrt{Dt}}\right]
\end{keyeq}
(рис.~\ref{fig:erf},\,\textit{в})\ispr{В оригинале (с.~433) в формуле
стоит $1-\operatorname{erf}$; такое решение отвечает обратному
начальному условию ($C_0$ при $x<0$).}. Каждому значению $C/C_0$
отвечает своё $z=x/2\sqrt{Dt}$: плоскость с любой постоянной
концентрацией удаляется от $x=0$ как $\sqrt{Dt}$, и только плоскость
$C=C_0/2$ ($z=0$) неподвижна.

\begin{figure}[ht]
\centering
\begin{subfigure}[b]{0.3\textwidth}\centering
\begin{tikzpicture}
\begin{axis}[width=1.02\textwidth, height=4.4cm, xmin=-3, xmax=3, ymin=0, ymax=1.15,
  xlabel={$x$}, ylabel={$C/C_0$}, xtick={0}, ytick={0,1}, axis lines=left]
  \addplot[fmBlue, thick] coordinates {(-3,0) (0,0) (0,1) (3,1)};
\end{axis}
\end{tikzpicture}
\caption{}
\end{subfigure}\hfill
\begin{subfigure}[b]{0.3\textwidth}\centering
\begin{tikzpicture}
\begin{axis}[width=1.15\textwidth, height=4.4cm, xmin=-3, xmax=3, ymin=0, ymax=1.15,
  xlabel={$x$}, xtick={0}, ytick={0,1}, axis lines=left]
  \addplot[fmGray, thin] table[x=x, y=g0] {figures/dif-erf-sum.dat};
  \addplot[fmGray, thin] table[x=x, y=g1] {figures/dif-erf-sum.dat};
  \addplot[fmGray, thin] table[x=x, y=g2] {figures/dif-erf-sum.dat};
  \addplot[fmGray, thin] table[x=x, y=g3] {figures/dif-erf-sum.dat};
  \addplot[fmGray, thin] table[x=x, y=g4] {figures/dif-erf-sum.dat};
  \addplot[fmGray, thin] table[x=x, y=g5] {figures/dif-erf-sum.dat};
  \addplot[fmRed, thick] table[x=x, y=sum] {figures/dif-erf-sum.dat};
\end{axis}
\end{tikzpicture}
\caption{}
\end{subfigure}\hfill
\begin{subfigure}[b]{0.3\textwidth}\centering
\begin{tikzpicture}
\begin{axis}[width=1.15\textwidth, height=4.4cm, xmin=-3, xmax=3, ymin=0, ymax=1.15,
  xlabel={$x$}, xtick={0}, ytick={0,0.5,1}, yticklabels={0,$\frac12$,1}, axis lines=left]
  \addplot[fmBlue, thick] table[x=x, y=L5] {figures/dif-erf.dat};
  \addplot[fmGreen!70!black, thick] table[x=x, y=L10] {figures/dif-erf.dat};
  \addplot[fmRed, thick] table[x=x, y=L20] {figures/dif-erf.dat};
  \draw[fmGray, dotted] (axis cs:-3,0.5) -- (axis cs:3,0.5);
\end{axis}
\end{tikzpicture}
\caption{}
\end{subfigure}
\caption{Диффузия между двумя полубесконечными телами (расчёт):
\textit{а}~— начальная «ступенька»; \textit{б}~— распределение как
сумма гауссиан от тонких слоёв (серые кривые~— вклады нескольких
слоёв); \textit{в}~— профили $\frac12[1+\operatorname{erf}(x/2\sqrt{Dt})]$
при $t_1$, $4t_1$, $16t_1$; все проходят через $C_0/2$ при $x=0$}
\label{fig:erf}
\end{figure}

Решение с функцией ошибок предполагает постоянный $D$, поэтому для
точного определения $D$ его применяют редко. Зато им удобно оценивать
количество вещества, проникшего в образец, когда $D$ меняется вдоль
зоны диффузии (например, при науглероживании или обезуглероживании
стали),~— получается среднее значение $D$.

\begin{primer}[Цементация]
Поверхность стали поддерживают при концентрации углерода $C_s$; в
глубине~— $C_\text{нач}\approx0$. Тогда
$C(x,t)=C_s\left[1-\operatorname{erf}\frac{x}{2\sqrt{Dt}}\right]$.
На какой глубине через 10~ч при $927$\,°C концентрация составит
половину поверхностной?
\begin{shagi}
\item $1-\operatorname{erf}z=\frac12$ при $z\approx0{,}48$.
\item $D$ углерода в аустените при 1200~К: $D=2{,}3\cdot10^{-5}\exp\bigl(-148\,000/(8{,}31\cdot1200)\bigr)\approx8\cdot10^{-12}$~м$^2$/с.
\item $x=2z\sqrt{Dt}=0{,}96\sqrt{8\cdot10^{-12}\cdot3{,}6\cdot10^4}\approx0{,}5$~мм.
\end{shagi}
Чтобы удвоить глубину слоя, цементацию пришлось бы вести 40~ч.
\end{primer}

\paragraph{Метод Больцмана--Матано}
При \emph{взаимной} диффузии двух металлов коэффициент $\widetilde D$
зависит от состава, и решение \eqref{eq:erf} неприменимо. Нужно решить
обратную задачу: по измеренному профилю $C(x)$ найти $\widetilde D(C)$.
\begin{shagi}
\item Подстановка Больцмана $\lambda=x/\sqrt t$ превращает уравнение
$\partial C/\partial t=\partial(\widetilde D\,\partial C/\partial x)/\partial x$
в обыкновенное: $-\frac\lambda2\frac{\d C}{\d\lambda}=\frac{\d}{\d\lambda}\left(\widetilde D\frac{\d C}{\d\lambda}\right)$.
\item Интегрируя от края зоны ($C'$, где $\d C/\d\lambda=0$) до $C$,
получаем $\widetilde D\frac{\d C}{\d\lambda}=-\frac12\int_{C'}^C\lambda\,\d C$.
\item Интегрирование от другого края ($C''$) должно дать то же самое.
Это возможно лишь при определённом выборе начала отсчёта~—
\term{плоскости Матано} $x_M$: $\int_{C'}^{C''}(x-x_M)\,\d C=0$.
\end{shagi}
\begin{keyeq}\label{eq:matano}
\widetilde D(C)=\frac1{2t}\left(\frac{\d x}{\d C}\right)_C\int_{C'}^C(x_M-x)\,\d C .
\end{keyeq}
Матано предложил вычислять всё графически (рис.~\ref{fig:matano}):
плоскость $x_M$ проводят так, чтобы заштрихованные площади слева и
справа от неё были равны; производную $\d x/\d C$ берут по наклону
кривой, интеграл~— как площадь. Самая трудоёмкая операция~— точное
определение производной.

% сгенерировано tools/gen_diffusion.py — не править вручную
\def\MatanoX{0.0043}

\begin{figure}[ht]
\centering
\begin{tikzpicture}
\begin{axis}[width=0.66\textwidth, height=5.4cm, xmin=-3, xmax=3, ymin=0, ymax=1.1,
  xlabel={$x$}, ylabel={$C$}, xtick={\MatanoX}, xticklabels={$x_M$}, ytick={0,1},
  yticklabels={$C'$,$C''$}, axis lines=left]
  \addplot[name path=prof, fmBlue, thick] table[x=x, y=C] {figures/dif-matano.dat};
  \path[name path=top] (axis cs:-3,1) -- (axis cs:3,1);
  \path[name path=bot] (axis cs:-3,0) -- (axis cs:3,0);
  \addplot[fmBlue!25] fill between[of=prof and bot, soft clip={domain=-3:\MatanoX}];
  \addplot[fmRed!25] fill between[of=prof and top, soft clip={domain=\MatanoX:3}];
  \draw[fmGray, dashed] (axis cs:\MatanoX,0) -- (axis cs:\MatanoX,1.05);
  \draw[fmGray, dotted] (axis cs:-3,1) -- (axis cs:3,1);
\end{axis}
\end{tikzpicture}
\caption{Профиль взаимной диффузии при $\widetilde D$, растущем с
концентрацией (расчёт): кривая несимметрична. Плоскость Матано $x_M$
выбрана так, что синяя и розовая площади равны}
\label{fig:matano}
\end{figure}
